% UET Thesis Pandoc LaTeX Template
% Customized for University of Engineering and Technology, VNU

\clearpage
\phantomsection

\setcounter{chapter}{1}
\chapter[Cơ sở toán học]{Cơ sở toán học}

Chương này trình bày các khái niệm toán học nền tảng được sử dụng trong
toàn bộ luận văn: cấu trúc trường hữu hạn, đường cong elliptic Short
Weierstrass và các phép toán nhóm trên nó
\citep{Silverman2009, Washington2008}, cùng với các kỹ thuật số học tối
ưu (Montgomery, Fermat) được cài đặt trực tiếp trong thư viện
\texttt{curve1024}.

\section{Định nghĩa Trường số nguyên
tố}\label{ux111ux1ecbnh-nghux129a-trux1b0ux1eddng-sux1ed1-nguyuxean-tux1ed1}

\textbf{Trường số nguyên tố} \(\mathbb{F}_p\) là tập
\(\{0, 1, \ldots, p-1\}\) với phép cộng và nhân lấy modulo \(p\), trong
đó \(p\) là số nguyên tố. Đây là trường hữu hạn nhỏ nhất có đặc số
\(p\).

Các tính chất quan trọng:

\begin{itemize}
\tightlist
\item
  \textbf{Đóng kín}: kết quả của mọi phép toán đều thuộc
  \(\mathbb{F}_p\).
\item
  \textbf{Nghịch đảo nhân}: mọi phần tử \(a \neq 0\) có nghịch đảo
  \(a^{-1}\) sao cho \(a \cdot a^{-1} \equiv 1 \pmod{p}\).
\item
  \textbf{Theo Định lý nhỏ Fermat}: \(a^{p-1} \equiv 1 \pmod{p}\) với
  mọi \(a \neq 0\), do đó \(a^{-1} = a^{p-2} \pmod{p}\).
\end{itemize}

\subsection{\texorpdfstring{Cài đặt: \texttt{PrimeFieldConfig} và
\texttt{PrimeFieldElement}}{Cài đặt: PrimeFieldConfig và PrimeFieldElement}}\label{cuxe0i-ux111ux1eb7t-primefieldconfig-vuxe0-primefieldelement}

Trong thư viện \texttt{curve1024}, trường \(\mathbb{F}_p\) được biểu
diễn qua trait \texttt{PrimeFieldConfig} (mang các hằng số tĩnh) và
struct \texttt{PrimeFieldElement\textless{}C\textgreater{}} (phần tử
trường):

\begin{Shaded}
\begin{Highlighting}[]
\KeywordTok{pub} \KeywordTok{trait}\NormalTok{ PrimeFieldConfig }\OperatorTok{\{}
    \KeywordTok{const}\NormalTok{ MODULUS}\OperatorTok{:}\NormalTok{ U1024}\OperatorTok{;}  \CommentTok{// p: số nguyên tố trường}
    \KeywordTok{const}\NormalTok{ R2}\OperatorTok{:}\NormalTok{      U1024}\OperatorTok{;}  \CommentTok{// R² mod p, dùng cho Montgomery}
    \KeywordTok{const}\NormalTok{ N\_PRIME}\OperatorTok{:}\NormalTok{ U1024}\OperatorTok{;}  \CommentTok{// {-}p\^{}\{{-}1\} mod 2\^{}1024, dùng cho REDC}
\OperatorTok{\}}

\KeywordTok{pub} \KeywordTok{struct}\NormalTok{ PrimeFieldElement}\OperatorTok{\textless{}}\NormalTok{C}\OperatorTok{:}\NormalTok{ PrimeFieldConfig}\OperatorTok{\textgreater{}} \OperatorTok{\{}
\NormalTok{    value}\OperatorTok{:}\NormalTok{ U1024}\OperatorTok{,}          \CommentTok{// biểu diễn Montgomery: a.R mod p}
\OperatorTok{\}}
\end{Highlighting}
\end{Shaded}

\captionof{lstlisting}{Định nghĩa PrimeFieldConfig và PrimeFieldElement}\label{lst:ch2_primefield}\vspace{12pt}

Thiết kế này cho phép trình biên dịch mở inline toàn bộ hằng số tại thời
gian biên dịch (zero-cost abstraction): hai trường khác nhau
\(\mathbb{F}_p\) và \(\mathbb{F}_r\) dùng chung code nhưng có tham số
riêng biệt không thể nhầm lẫn.

\subsection{Các phép toán cơ
bản}\label{cuxe1c-phuxe9p-touxe1n-cux1a1-bux1ea3n}

\textbf{Cộng và trừ} thực hiện bằng cộng/trừ số nguyên rồi hiệu chỉnh
modulo:

\begin{Shaded}
\begin{Highlighting}[]
\CommentTok{// Cộng: (a + b) mod p}
\KeywordTok{fn}\NormalTok{ add(}\KeywordTok{self}\OperatorTok{,}\NormalTok{ rhs}\OperatorTok{:} \DataTypeTok{Self}\NormalTok{) }\OperatorTok{{-}\textgreater{}} \DataTypeTok{Self} \OperatorTok{\{}
    \KeywordTok{let}\NormalTok{ (sum}\OperatorTok{,}\NormalTok{ carry) }\OperatorTok{=} \KeywordTok{self}\OperatorTok{.}\NormalTok{value}\OperatorTok{.}\NormalTok{carrying\_add(}\OperatorTok{\&}\NormalTok{rhs}\OperatorTok{.}\NormalTok{value)}\OperatorTok{;}
    \KeywordTok{let}\NormalTok{ (sub}\OperatorTok{,}\NormalTok{ borrow) }\OperatorTok{=}\NormalTok{ sum}\OperatorTok{.}\NormalTok{borrowing\_sub(}\OperatorTok{\&}\PreprocessorTok{C::}\NormalTok{MODULUS)}\OperatorTok{;}
    \CommentTok{// Chọn kết quả không cần phép chia}
\NormalTok{    conditional\_select(sum}\OperatorTok{,}\NormalTok{ sub}\OperatorTok{,}\NormalTok{ carry }\OperatorTok{||} \OperatorTok{!}\NormalTok{borrow)}
\OperatorTok{\}}
\end{Highlighting}
\end{Shaded}

\captionof{lstlisting}{Cài đặt phép cộng trên trường hữu hạn modulo $p$}\label{lst:ch2_add}\vspace{12pt}

Phép cộng không cần phép chia đầy đủ --- chỉ cần một phép trừ có điều
kiện. Chi phí: \(O(n)\) với \(n\) là số limb (16 limb x 64 bit = 1024
bit).

\textbf{Nhân} dùng phép nhân Montgomery (xem §2.4).

\textbf{Nghịch đảo} dùng lũy thừa Fermat: \(a^{-1} = a^{p-2} \pmod{p}\),
thực hiện bởi hàm \texttt{pow} với phép nhân bình phương liên tiếp
(square-and-multiply):

\begin{Shaded}
\begin{Highlighting}[]
\KeywordTok{pub} \KeywordTok{fn}\NormalTok{ inv(}\OperatorTok{\&}\KeywordTok{self}\NormalTok{) }\OperatorTok{{-}\textgreater{}} \DataTypeTok{Self} \OperatorTok{\{}
    \KeywordTok{self}\OperatorTok{.}\NormalTok{pow(}\PreprocessorTok{C::}\NormalTok{MODULUS }\OperatorTok{{-}} \DecValTok{2}\NormalTok{)  }\CommentTok{// a\^{}(p{-}2) mod p}
\OperatorTok{\}}
\end{Highlighting}
\end{Shaded}

\captionof{lstlisting}{Cài đặt phép nghịch đảo bằng Định lý nhỏ Fermat}\label{lst:ch2_inv}\vspace{12pt}

\section{Đường cong elliptic Short
Weierstrass}\label{ux111ux1b0ux1eddng-cong-elliptic-short-weierstrass}

\subsection{Định nghĩa}\label{ux111ux1ecbnh-nghux129a}

Trên \(\mathbb{F}_p\) với \(p > 3\), \textbf{đường cong elliptic dạng
Short Weierstrass} là tập hợp:

\[E(\mathbb{F}_p) = \{(x, y) \in \mathbb{F}_p \times \mathbb{F}_p \mid y^2 = x^3 + ax + b\} \cup \{\mathcal{O}\}\]

trong đó \(\mathcal{O}\) là \textbf{điểm vô cực} (point at infinity) ---
phần tử đặc biệt không thuộc mặt phẳng affine, đóng vai trò phần tử đơn
vị của nhóm.

\textbf{Điều kiện không suy biến:} Biệt thức
\(\Delta = -16(4a^3 + 27b^2) \neq 0\), đảm bảo đường cong không có điểm
kỳ dị (cusp hay node). Kiểm tra này được thực hiện ngầm qua quá trình
CM: với \(D = -3\), luôn có \(a = 0\) và điều kiện thu gọn thành
\(b \neq 0\).

\subsection{Cấu trúc nhóm}\label{cux1ea5u-truxfac-nhuxf3m}

Tập \(E(\mathbb{F}_p)\) tạo thành một \textbf{nhóm Abel hữu hạn} với
phép cộng điểm. Bậc của nhóm \(\#E(\mathbb{F}_p) = p + 1 - t\) (định lý
Hasse \citep{Silverman2009}), trong đó \(|t| \leq 2\sqrt{p}\) là vết
Frobenius.

Trong cài đặt:

\begin{Shaded}
\begin{Highlighting}[]
\KeywordTok{pub} \KeywordTok{struct}\NormalTok{ AffinePoint}\OperatorTok{\textless{}}\NormalTok{C}\OperatorTok{:}\NormalTok{ SWCurveConfig}\OperatorTok{\textgreater{}} \OperatorTok{\{}
    \KeywordTok{pub}\NormalTok{ x}\OperatorTok{:}\NormalTok{           FieldElement}\OperatorTok{\textless{}}\PreprocessorTok{C::}\NormalTok{BaseField}\OperatorTok{\textgreater{},}
    \KeywordTok{pub}\NormalTok{ y}\OperatorTok{:}\NormalTok{           FieldElement}\OperatorTok{\textless{}}\PreprocessorTok{C::}\NormalTok{BaseField}\OperatorTok{\textgreater{},}
    \KeywordTok{pub}\NormalTok{ is\_infinite}\OperatorTok{:} \DataTypeTok{bool}\OperatorTok{,}
\OperatorTok{\}}
\end{Highlighting}
\end{Shaded}

\captionof{lstlisting}{Định nghĩa cấu trúc điểm AffinePoint}\label{lst:ch2_affine}\vspace{12pt}

\section{Phép cộng điểm (Point
Addition)}\label{phuxe9p-cux1ed9ng-ux111iux1ec3m-point-addition}

\subsection{Quy tắc hình học}\label{quy-tux1eafc-huxecnh-hux1ecdc}

Phép cộng trên \(E\) được định nghĩa bằng quy tắc ``dây cung và tiếp
tuyến'':

\begin{itemize}
\tightlist
\item
  \textbf{Phần tử đơn vị:} \(P + \mathcal{O} = \mathcal{O} + P = P\) với
  mọi \(P\).
\item
  \textbf{Phần tử nghịch đảo:} \(-( x, y) = (x, -y)\), vì
  \(P + (-P) = \mathcal{O}\).
\item
  \textbf{Cộng hai điểm phân biệt} \(P_1 \neq \pm P_2\): vẽ đường thẳng
  qua \(P_1, P_2\), điểm giao thứ ba với đường cong là \(R'\), phản
  chiếu qua trục \(x\) cho \(P_3 = P_1 + P_2 = -R'\).
\item
  \textbf{Nhân đôi} \(P + P = 2P\): dùng đường tiếp tuyến tại \(P\).
\end{itemize}

\subsection{Công thức đại
số}\label{cuxf4ng-thux1ee9c-ux111ux1ea1i-sux1ed1}

Cho \(P_1 = (x_1, y_1)\) và \(P_2 = (x_2, y_2)\) trên
\(y^2 = x^3 + ax + b\), điểm tổng \(P_3 = (x_3, y_3)\):

\textbf{Cộng hai điểm phân biệt} (\(P_1 \neq P_2\)):
\[\lambda = \frac{y_2 - y_1}{x_2 - x_1}, \quad x_3 = \lambda^2 - x_1 - x_2, \quad y_3 = \lambda(x_1 - x_3) - y_1\]

\textbf{Nhân đôi} (\(P_1 = P_2\)):
\[\lambda = \frac{3x_1^2 + a}{2y_1}, \quad x_3 = \lambda^2 - 2x_1, \quad y_3 = \lambda(x_1 - x_3) - y_1\]

Mỗi công thức cần \textbf{1 phép nghịch đảo} (tính \(\lambda\)), 2--3
phép nhân, và 4--6 phép cộng/trừ trong \(\mathbb{F}_p\).

\subsection{Cài đặt}\label{cuxe0i-ux111ux1eb7t}

Các bước cộng và nhân đôi điểm trên đường cong được cài đặt trực tiếp từ
công thức đại số, thông qua hai hàm \texttt{add()} và \texttt{double()}.
Đặc biệt, hệ thống tích hợp sẵn hàm \texttt{is\_on\_curve()} trong quá
trình khởi tạo điểm kết quả thông qua hàm \texttt{new()} để kiểm tra và
loại bỏ ngay lập tức các điểm không hợp lệ. Điều này giúp hệ thống bắt
lỗi sớm và miễn nhiễm hoàn toàn với tấn công đường cong không hợp lệ
(Invalid Curve Attack).

Chi tiết mã nguồn cài đặt các hàm \texttt{add()} và \texttt{double()}
được trình bày tại mã nguồn công khai của dự án \citep{LumenMath}.

\section{Phép nhân vô hướng (Scalar
Multiplication)}\label{phuxe9p-nhuxe2n-vuxf4-hux1b0ux1edbng-scalar-multiplication}

\subsection{Thuật toán
Double-and-Add}\label{thuux1eadt-touxe1n-double-and-add}

Phép nhân vô hướng \([k]P = P + P + \cdots + P\) (\(k\) lần) được tính
hiệu quả bằng \textbf{Double-and-Add} (tương tự binary exponentiation),
với độ phức tạp \(O(\log k)\) thay vì \(O(k)\):

\begin{verbatim}
Input: P in E(F_p), k in [0, r-1]
Output: [k]P

R ← O (điểm vô cực)
B ← P
for i = 0 to 1023:
    if bit i của k bằng 1:
        R ← R + B
    B ← 2B
return R
\end{verbatim}

Với \(k \approx 2^{512}\) (1024 bit), thuật toán thực hiện tối đa 1024
lần nhân đôi và trung bình 512 lần cộng điểm.

\subsection{Cài đặt}\label{cuxe0i-ux111ux1eb7t-1}

\begin{Shaded}
\begin{Highlighting}[]
\KeywordTok{pub} \KeywordTok{fn}\NormalTok{ mul(}\OperatorTok{\&}\KeywordTok{self}\OperatorTok{,}\NormalTok{ scalar}\OperatorTok{:} \OperatorTok{\&}\NormalTok{U1024) }\OperatorTok{{-}\textgreater{}} \DataTypeTok{Self} \OperatorTok{\{}
    \KeywordTok{let} \KeywordTok{mut}\NormalTok{ result }\OperatorTok{=} \DataTypeTok{Self}\PreprocessorTok{::}\NormalTok{infinite()}\OperatorTok{;}
    \KeywordTok{let} \KeywordTok{mut}\NormalTok{ base   }\OperatorTok{=} \OperatorTok{*}\KeywordTok{self}\OperatorTok{;}
    \ControlFlowTok{for}\NormalTok{ i }\KeywordTok{in} \DecValTok{0}\OperatorTok{..}\DecValTok{1024} \OperatorTok{\{}
        \KeywordTok{let}\NormalTok{ limb\_idx }\OperatorTok{=}\NormalTok{ i }\OperatorTok{/} \DecValTok{64}\OperatorTok{;}
        \KeywordTok{let}\NormalTok{ bit\_idx  }\OperatorTok{=}\NormalTok{ i }\OperatorTok{\%} \DecValTok{64}\OperatorTok{;}
        \ControlFlowTok{if}\NormalTok{ (scalar}\OperatorTok{.}\DecValTok{0}\NormalTok{[limb\_idx] }\OperatorTok{\textgreater{}\textgreater{}}\NormalTok{ bit\_idx) }\OperatorTok{\&} \DecValTok{1} \OperatorTok{==} \DecValTok{1} \OperatorTok{\{}
\NormalTok{            result }\OperatorTok{=}\NormalTok{ result}\OperatorTok{.}\NormalTok{add(}\OperatorTok{\&}\NormalTok{base)}\OperatorTok{;}
        \OperatorTok{\}}
\NormalTok{        base }\OperatorTok{=}\NormalTok{ base}\OperatorTok{.}\NormalTok{double()}\OperatorTok{;}
    \OperatorTok{\}}
\NormalTok{    result}
\OperatorTok{\}}
\end{Highlighting}
\end{Shaded}

\captionof{lstlisting}{Thuật toán nhân vô hướng Double-and-Add}\label{lst:ch2_mul}\vspace{12pt}

Mỗi lần gọi \texttt{mul} thực hiện đúng 1024 lần \texttt{double} và tối
đa 1024 lần \texttt{add}. Chi phí tính toán chủ yếu đến từ phép nghịch
đảo trong \(\mathbb{F}_p\) (mỗi lần \texttt{add}/\texttt{double} cần 1
nghịch đảo, tức \(p^{p-2}\) --- một lũy thừa 1024-bit).

\subsection{Tầm quan trọng}\label{tux1ea7m-quan-trux1ecdng}

Phép nhân vô hướng là hàm cốt lõi của mọi sơ đồ mật mã ECC:

\begin{longtable}[]{@{}
  >{\raggedright\arraybackslash}p{(\columnwidth - 4\tabcolsep) * \real{0.4062}}
  >{\raggedright\arraybackslash}p{(\columnwidth - 4\tabcolsep) * \real{0.2656}}
  >{\raggedright\arraybackslash}p{(\columnwidth - 4\tabcolsep) * \real{0.3281}}@{}}
\caption{Vai trò của phép nhân vô hướng trong các sơ đồ mật mã
ECC}\tabularnewline
\toprule\noalign{}
\begin{minipage}[b]{\linewidth}\raggedright
Thao tác
\end{minipage} & \begin{minipage}[b]{\linewidth}\raggedright
Ký hiệu
\end{minipage} & \begin{minipage}[b]{\linewidth}\raggedright
Ứng dụng
\end{minipage} \\
\midrule\noalign{}
\endfirsthead
\toprule\noalign{}
\begin{minipage}[b]{\linewidth}\raggedright
Thao tác
\end{minipage} & \begin{minipage}[b]{\linewidth}\raggedright
Ký hiệu
\end{minipage} & \begin{minipage}[b]{\linewidth}\raggedright
Ứng dụng
\end{minipage} \\
\midrule\noalign{}
\endhead
\bottomrule\noalign{}
\endlastfoot
Sinh khóa công khai & \(Q = [d]G\) & Từ khóa riêng \(d\) \\
Cam kết ký (Schnorr/ECDSA) & \(R = [k]G\) & Nonce commitment \\
Xác minh Schnorr & \([s]G\) và \([e]Q\) & Hai lần nhân vô hướng \\
Xác minh ECDSA & \([u_1]G + [u_2]Q\) & Hai lần nhân vô hướng \\
\end{longtable}

\section{Tối ưu số học: Phép nhân
Montgomery}\label{tux1ed1i-ux1b0u-sux1ed1-hux1ecdc-phuxe9p-nhuxe2n-montgomery}

\subsection{Vấn đề}\label{vux1ea5n-ux111ux1ec1}

Phép nhân thông thường \(a \cdot b \pmod{p}\) đòi hỏi một phép chia đầy
đủ để lấy số dư --- rất đắt với số 1024-bit. Montgomery đề xuất thực
hiện phép nhân trong \textbf{không gian Montgomery} để tránh phép chia
\citep{Montgomery1985}.

\subsection{Biến đổi
Montgomery}\label{biux1ebfn-ux111ux1ed5i-montgomery}

Chọn \(R = 2^{1024}\) (cơ số nguyên tố cùng nhau với \(p\)).
\textbf{Dạng Montgomery} của \(a\) là \(\hat{a} = a \cdot R \pmod{p}\).

Phép nhân Montgomery (\textbf{REDC}) tính
\(\hat{a} \cdot \hat{b} \cdot R^{-1} \pmod{p}\) từ \(\hat{a}\) và
\(\hat{b}\), chỉ dùng phép chia cho \(R\) (shift phải 1024 bit --- rất
rẻ) thay vì chia cho \(p\).

\begin{longtable}[]{@{}lll@{}}
\caption{So sánh chi phí phép nhân thông thường và
Montgomery}\tabularnewline
\toprule\noalign{}
Phép tính & Thông thường & Montgomery \\
\midrule\noalign{}
\endfirsthead
\toprule\noalign{}
Phép tính & Thông thường & Montgomery \\
\midrule\noalign{}
\endhead
\bottomrule\noalign{}
\endlastfoot
\(a \cdot b \pmod{p}\) & cần chia cho \(p\) & shift + trừ có điều
kiện \\
Chi phí & \(O(n^2)\) + chia & \(O(n^2)\) không chia \\
\end{longtable}

\textbf{Quy trình REDC} (với ba hằng số biên dịch: \(R^2 \bmod p\),
\(-p^{-1} \bmod R\)):

\begin{enumerate}
\def\labelenumi{\arabic{enumi}.}
\tightlist
\item
  Nhân \(\text{lo}, \text{hi} = \hat{a} \cdot \hat{b}\)
\item
  \(m = \text{lo} \cdot N' \bmod R\) với \(N' = -p^{-1} \bmod R\)
\item
  \(t = (\text{hi}, \text{lo}) + m \cdot p\) (chỉ giữ phần trên \(R\))
\item
  Nếu \(t \geq p\): \(t \leftarrow t - p\)
\end{enumerate}

\begin{Shaded}
\begin{Highlighting}[]
\KeywordTok{fn}\NormalTok{ reduce(lo}\OperatorTok{:} \OperatorTok{\&}\NormalTok{U1024}\OperatorTok{,}\NormalTok{ hi}\OperatorTok{:} \OperatorTok{\&}\NormalTok{U1024) }\OperatorTok{{-}\textgreater{}}\NormalTok{ U1024 }\OperatorTok{\{}
    \KeywordTok{let}\NormalTok{ (m}\OperatorTok{,}\NormalTok{ \_)       }\OperatorTok{=}\NormalTok{ lo}\OperatorTok{.}\NormalTok{widening\_mul(}\OperatorTok{\&}\PreprocessorTok{C::}\NormalTok{N\_PRIME)}\OperatorTok{;}   \CommentTok{// m = lo . N\textquotesingle{} mod R}
    \KeywordTok{let}\NormalTok{ (mn\_lo}\OperatorTok{,}\NormalTok{ mn\_hi) }\OperatorTok{=}\NormalTok{ m}\OperatorTok{.}\NormalTok{widening\_mul(}\OperatorTok{\&}\PreprocessorTok{C::}\NormalTok{MODULUS)}\OperatorTok{;}  \CommentTok{// m . p}
    \KeywordTok{let}\NormalTok{ (\_}\OperatorTok{,}\NormalTok{ carry\_lo)  }\OperatorTok{=}\NormalTok{ lo}\OperatorTok{.}\NormalTok{carrying\_add(}\OperatorTok{\&}\NormalTok{mn\_lo)}\OperatorTok{;}
    \KeywordTok{let}\NormalTok{ (t}\OperatorTok{,}\NormalTok{ carry\_hi)  }\OperatorTok{=}\NormalTok{ hi}\OperatorTok{.}\NormalTok{carrying\_add(}\OperatorTok{\&}\NormalTok{mn\_hi)}\OperatorTok{;}
    \CommentTok{// Bước hiệu chỉnh cuối: trừ p nếu cần}
    \KeywordTok{let}\NormalTok{ (sub}\OperatorTok{,}\NormalTok{ borrow) }\OperatorTok{=}\NormalTok{ t}\OperatorTok{.}\NormalTok{borrowing\_sub(}\OperatorTok{\&}\PreprocessorTok{C::}\NormalTok{MODULUS)}\OperatorTok{;}
    \ControlFlowTok{if}\NormalTok{ carry\_hi }\OperatorTok{||}\NormalTok{ carry\_lo }\OperatorTok{\{}\NormalTok{ sub }\OperatorTok{\}} \ControlFlowTok{else} \ControlFlowTok{if} \OperatorTok{!}\NormalTok{borrow }\OperatorTok{\{}\NormalTok{ sub }\OperatorTok{\}} \ControlFlowTok{else} \OperatorTok{\{}\NormalTok{ t }\OperatorTok{\}}
\OperatorTok{\}}
\end{Highlighting}
\end{Shaded}

\captionof{lstlisting}{Hàm giảm bậc Montgomery (REDC)}\label{lst:ch2_reduce}\vspace{12pt}

\subsection{Chuyển đổi vào/ra không gian
Montgomery}\label{chuyux1ec3n-ux111ux1ed5i-vuxe0ora-khuxf4ng-gian-montgomery}

\begin{itemize}
\tightlist
\item
  \textbf{Chuyển vào:} \texttt{PrimeFieldElement::new(a)} tính
  \(a \cdot R^2 \cdot R^{-1} = a \cdot R \pmod{p}\).
\item
  \textbf{Chuyển ra:} \texttt{to\_u1024()} gọi
  \texttt{reduce(value,\ 0)} = \(\hat{a} \cdot R^{-1} = a \pmod{p}\).
\end{itemize}

Trong thực tế, hầu hết tính toán diễn ra trong không gian Montgomery ---
chỉ chuyển ra khi cần so sánh hay xuất kết quả, giúp loại bỏ hoàn toàn
phép chia trong vòng lặp tính toán chính.

\section{Lũy thừa nhanh
(Square-and-Multiply)}\label{lux169y-thux1eeba-nhanh-square-and-multiply}

Hàm \texttt{pow} thực hiện lũy thừa \(a^e \pmod{p}\) bằng phương pháp
\textbf{square-and-multiply} (nhân bình phương liên tiếp). Việc cài đặt
cấp thấp tiềm ẩn rủi ro lộ lọt thông tin khóa qua độ chênh lệch thời
gian thực thi của các khối lệnh điều kiện.

Để bảo vệ hệ thống trước tấn công kênh kề (timing side-channel), thay vì
sử dụng cấu trúc rẽ nhánh \texttt{if-else} (mở ra nguy cơ phán đoán bit
số mũ dựa trên thời gian thực thi của CPU branch-prediction), thủ tục sử
dụng cách tính chọn hằng-thời-gian (constant-time execution):

\begin{Shaded}
\begin{Highlighting}[]
\CommentTok{// Cả hai nhánh điều kiện đều được tính toán}
\KeywordTok{let}\NormalTok{ product }\OperatorTok{=}\NormalTok{ res }\OperatorTok{*}\NormalTok{ base}\OperatorTok{;}
\CommentTok{// Chỉ định phép gán qua hàm chuyển đổi bit, không dùng rẽ nhánh điều khiển}
\NormalTok{res }\OperatorTok{=} \DataTypeTok{Self}\PreprocessorTok{::}\NormalTok{conditional\_select(}\OperatorTok{\&}\NormalTok{res}\OperatorTok{,} \OperatorTok{\&}\NormalTok{product}\OperatorTok{,}\NormalTok{ bit}\OperatorTok{.}\NormalTok{into())}\OperatorTok{;}
\end{Highlighting}
\end{Shaded}

\captionof{lstlisting}{Phép gán hằng-thời-gian trong lũy thừa chống tấn công kênh kề}\label{lst:ch2_pow_constant_time}\vspace{12pt}

Hàm \texttt{conditional\_select} đảm bảo cả hai nhánh (khi bit bằng 0
hoặc 1) đều tiêu tốn một lượng chu kỳ lệnh (clock cycles) CPU bằng nhau
hoàn toàn. Chi tiết mã nguồn vòng lặp 1024-bit của hàm \texttt{pow} được
đính kèm tại mã nguồn công khai của dự án \citep{LumenMath}.

Nghịch đảo \(a^{-1} = a^{p-2}\) dùng \texttt{pow} với \(e = p - 2\):
1024 lần bình phương và \textasciitilde512 lần nhân (trung bình), tổng
\(\approx 1536\) phép nhân Montgomery.

\FloatBarrier
